\section{A general start}\label{sec:general-start}

Kagey's ball picks its first direction uniformly. Suppose instead that it
starts in state $a$ with probability $\mu_a$, and moves as before. The start
only affects the first run, so one might expect it to matter little. It does
matter: it changes the constants $a_k$ by an amount of order $1$, which is
exactly the order at which the faces are separated. So a start can decide which
face wins.

For a count vector $n$ with support $F$, Proposition~\ref{prop:dpmf}(b) gives
\[
 P^\mu_N(n)=p^{N-1}\sum_{a\in F}\mu_aS^{(a)}_n(u)
 =p^{N-1}S_n(u)\sum_{a\in F}\mu_a\omega_a(n),\qquad
 \omega_a(n)=\frac{S^{(a)}_n(u)}{S_n(u)}.
\]
So the start enters only through the weights $\omega_a(n)$, which sum to $1$. A
vertex has $V^\mu_i=\mu_ip^{N-1}$. Throughout this section, $F$ is a face with
$k=|F|\ge2$ states and $\mu(F)>0$ (if $\mu(F)=0$, every bin with support $F$ has
probability $0$), and $i$ is a vertex with $\mu_i>0$.

\begin{lemma}[start weights]\label{lem:start-weights}
Let $n$ be a positive count vector on $F$.
\begin{enumerate}[(1)]
\item If $n_a=n_b$, then $\omega_a(n)=\omega_b(n)$. So $\omega_a(n)=1/k$ for
all $a\in F$ when $n$ is balanced and $k$ divides $N$.
\item Fix $C>0$. There are $C_1$ and $N_0$, depending only on $d$, $k$ and $C$,
such that for $N\ge N_0$, $0<T\le CL$ and every balanced bin $n$ of $F$,
\[
 \Bigl|\sum_{a\in F}\mu_a\omega_a(n)-\frac{\mu(F)}k\Bigr|
 \le\mu(F)\,C_1\frac{(1+T)^2}N.
\]
\end{enumerate}
\end{lemma}

In words, a balanced bin forgets where the path started, up to a relative error
$O(T^2/N)$. The proof is in Appendix~\ref{app:start-field}. For bins that are
not balanced the weights are not uniform, and the next remark says what they
are.

\begin{remark}[where a conditioned path starts]\label{rem:start-weights}
Under the uniform start, $\omega_a(n)$ is the probability that the chain starts
in $a$ given $K_N=n$. Fix a compact set $\mathcal K$ in the open face and
$0<c<C$, and put $\alpha=n/N$ and $A=\sum_{b\in F}\sqrt{\alpha_b}$. Then
$\omega_a(n)=\sqrt{\alpha_a}/A+O(1/L)$, uniformly for $\alpha\in\mathcal K$ and
$cL\le T\le CL$ (proof omitted; checked numerically, see
Section~\ref{sec:verify}). So a path conditioned to have fractions $\alpha$
starts in $a$ with probability close to $\sqrt{\alpha_a}/A$, and not
$\alpha_a$. In the notation of Proposition~\ref{prop:kirchhoff}, $\sqrt\alpha$
is the Perron vector $h$ of the tilted generator.
\end{remark}

\begin{theorem}[face-center crossings with a start]\label{thm:start-crossings}
Let $n$ be a balanced bin of $F$.
\begin{enumerate}[(1)]
\item For every $N\ge k$, the ratio $P^\mu_N(n)/V^\mu_i$ is a polynomial in $u$
with nonnegative coefficients, zero constant term and lowest term
$(\mu(F)/\mu_i)(k-1)!\,u^{k-1}$. So it equals $1$ at exactly one $p\in(0,1)$, and
the bin is more likely than the vertex below this $p$. If $k$ divides $N$, then
$P^\mu_N(n)=\frac{\mu(F)}k\,p^{N-1}S_n(u)=\frac{d\mu(F)}k\,P_N(n)$.
So the crossing solves $S_n(u)=k\mu_i/\mu(F)$, its $u$ decreases strictly along
$N\in k\Z$, and at $N=k$ it is $u=(\mu_i/(\mu(F)(k-1)!))^{1/(k-1)}$.
\item Fix $C>0$ and let $N_0$ be as in Lemma~\ref{lem:start-weights}(2). For
$N\ge N_0$, uniformly for $0<T\le CL$ and over all $\mu$,
$P^\mu_N(n)=\frac{\mu(F)}k\,p^{N-1}S_n(u)\bigl(1+O((1+T)^2/N)\bigr)$.
\item Put $a_k(\mu)=a_k+\frac1{k-1}\log\frac{k\mu_i}{\mu(F)}$. Fix $0<c<C$. For
large $N$, uniformly for $cL\le T\le CL$ and over all $\mu$,
\[
 \log\frac{P^\mu_N(n)}{V^\mu_i}=(k-1)\bigl[\Psi_N(T)-a_k(\mu)\bigr]
 -\frac{k-1}{8T}+O\Bigl(\frac1{T^2}+\frac{T^2}N\Bigr).
\]
So the value $z=z^\mu_N$ of $T$ at the crossing satisfies
$z+\frac12\log z=L+a_k(\mu)+\frac1{8z}+O(L^{-2})$ and the expansion
\eqref{eq:simplex-root-expansion} with $a_k(\mu)$ in place of $a_k$. With
$T=L-\frac12\log L+t$ the ratio tends to $e^{(k-1)(t-a_k(\mu))}$, uniformly on
compact sets of $t$.
\end{enumerate}
\end{theorem}

In words, a start changes only the constant of each crossing. It lowers $a_k$
by $\frac1{k-1}\log(\mu(F)/(k\mu_i))$, the logarithm of the average start
weight on the face against the start weight of the vertex, shared among the
$k-1$ new states. The error terms do not depend on $\mu$.

\begin{proof}
(1) Each $S^{(a)}_n$ has nonnegative coefficients, zero constant term and lowest
term $(k-1)!\,u^{k-1}$, from the $(k-1)!$ orders of the states that start with
$a$, with each state in a single run. So the ratio increases strictly from $0$
to $\infty$ in $u$, and $u=r/p$ decreases from $\infty$ to $0$ on $(0,1)$. If
$k$ divides $N$, Lemma~\ref{lem:start-weights}(1) gives
$\sum_a\mu_a\omega_a(n)=\mu(F)/k$. The rest of (1) is about $S_n$, which
increases strictly along $N\in k\Z$ at every $u>0$ and equals $k!\,u^{k-1}$ at
$N=k$, by the proof of the exact part of Theorem~\ref{thm:simplex-crossing}.

(2) is Lemma~\ref{lem:start-weights}(2), divided by $\mu(F)/k$.

(3) By (2), $\log(P^\mu_N(n)/V^\mu_i)=\log S_n(u)-\log(k\mu_i/\mu(F))+O(T^2/N)$,
and $S_n(u)=C_{N,k}/V_N$. So Theorem~\ref{thm:simplex-window}(i) gives the
display. The rest follows as in the proof of Theorem~\ref{thm:simplex-crossing},
since only the constant has changed.
\end{proof}

In the notation of Paper~C~\cite{paperC}, where $K_F$ is the second-order
constant of a face, $K_F=\mu(F)k^{k-1/2}(4\pi)^{-(k-1)/2}$ and
$K_{\{i\}}=\mu_i$ here, so $a_k(\mu)=-\frac1{k-1}\log(K_F/K_{\{i\}})$. So $K_F$
depends on the start only through $\mu(F)$, and the uniform start gives
$a_k(\mu)=a_k$.

If $\mu$ is constant on $F$, say $\mu_a=m$ for $a\in F$, then
$P^\mu_N(n)=m\,p^{N-1}S_n(u)$ for every $n$ with support in $F$. So the
statements of Sections~\ref{sec:simplex-thresholds}
and~\ref{sec:frontier-boundary} about such bins hold verbatim, with $V_N$
replaced by $V^\mu_a=m\,p^{N-1}$ for a vertex $a\in F$. In particular
Theorem~\ref{thm:simplex-balance} holds on $F$. Otherwise balancing can fail.

\begin{example}[balancing can fail]\label{ex:start-balance}
Let $d=2$, $N=4$ and $\mu=(\frac12+\varepsilon,\frac12-\varepsilon)$ with
$0<\varepsilon\le\frac12$. The words with counts $(3,1)$ are $1112$, $2111$,
$1121$ and $1211$, and those with counts $(2,2)$ are $1122$, $2211$, $1221$,
$2112$, $1212$ and $2121$. Counting changes gives
$P^\mu_4(3,1)=p^3(u+2\mu_1u^2)$ and $P^\mu_4(2,2)=p^3(u+u^2+u^3)$, so
\[
 P^\mu_4(3,1)-P^\mu_4(2,2)=p^3u^2(2\varepsilon-u).
\]
So the unbalanced bin is more likely when $u<2\varepsilon$, and for
$\mu=\delta_1$ at every $\frac12<p<1$.
\end{example}

This example only shows that balancing can fail. For every $\mu$ that is not
constant on $F$, the next-term expansion that Paper~C states without the
details~\cite{paperC} predicts more. For all large $N$ and $cL\le T\le CL$, it
predicts that the largest bin on $F$ beats the balanced bins by a factor
$1+\frac1{4\mu(F)^2T}\sum_{a\in F}(\mu_a-\mu(F)/k)^2+o(1/T)$. So balancing should
fail there. We do not prove this.

Now we compare all faces at once. Each face gets a line in the window
coordinate $t$, and the highest line wins.

\begin{remark}[global modes with a start]\label{rem:start-modes}
This remark uses a result of Paper~C. Order the start weights as
$\mu_{(1)}\ge\cdots\ge\mu_{(d)}$, put $M_k=\mu_{(1)}+\cdots+\mu_{(k)}$ and
$T=L-\frac12\log L+t$, and define the lines
\[
 \ell^\mu_1(t)=\log\mu_{(1)},\qquad
 \ell^\mu_k(t)=\log\frac{M_k}k+(k-1)(t-a_k)\quad(2\le k\le d).
\]
By Theorem~\ref{thm:start-crossings}(3), a balanced bin $n$ of a face $S$ with
$k\ge2$ states has $\log(P^\mu_N(n)/p^{N-1})=\log\frac{\mu(S)}k+(k-1)(t-a_k)+o(1)$,
and a vertex $i$ has $\log\mu_i$. Paper~C proves, from its general theorems on
face maxima and tie windows with the constants $K_F$ above, that the best bin of
each face has the same limit and that the lines decide the modes. So if
$\max_k\ell^\mu_k(t)$ is attained only at $k=k^*$, then for this $t$ and all
large $N$ every global mode of $P^\mu_N$ has exactly $k^*$ positive coordinates,
and its support $S$ has $\mu(S)=M_{k^*}$.

The winning size is nondecreasing in $t$. Indeed, if $k_1$ wins at $t_1$ and
$k_2$ at $t_2>t_1$, adding the 2 comparisons gives $(k_2-k_1)(t_2-t_1)\ge0$. Put
$t^*=a_d+\frac1{d-1}\log(d\mu_{(1)})$, where $\ell^\mu_1$ and $\ell^\mu_d$
cross. Their maximum is convex with its only kink at $t^*$, and the other lines
have slopes strictly between $0$ and $d-1$. So some size $2\le k\le d-1$ is the
unique winner on an open interval of $t$ if and only if
$\ell^\mu_k(t^*)>\ell^\mu_1(t^*)$ for some such $k$. For the uniform start
$\ell^\mu_k(t^*)-\ell^\mu_1(t^*)=(k-1)(a_d-a_k)<0$, which agrees with
Corollary~\ref{cor:simplex-modes}.
\end{remark}

\begin{remark}[the start decides which face wins]\label{ex:K3-start}
Let $d=3$ and $\mu=(\frac12,\frac12,0)$, and put
$t_{\rm ef}=\log\frac32+2a_3-a_2=\frac12\log\frac{32\pi}{243}$. In the terms of
Remark~\ref{rem:start-modes}, the lines are $\log\frac12$ for the vertices $1$
and $2$, $\log\frac12+t-a_2$ for the edge $\{1,2\}$, $\log\frac14+t-a_2$ for the
edges through $3$, and $\log\frac13+2(t-a_3)$ for the full face. The edge
$\{1,2\}$ has the highest line exactly when $a_2<t<t_{\rm ef}$. So for every
fixed $t$ in this window and all large $N$, the global modes are the balanced
bins of the edge $\{1,2\}$, since Theorem~\ref{thm:simplex-balance} applies on
the edge, where $\mu$ is constant. For fixed $t<a_2$ the modes are $(N,0,0)$ and
$(0,N,0)$, and for fixed $t>t_{\rm ef}$ they have 3 positive counts. The window
is $(-0.4673558,-0.4412978)$, of width $\log\frac{16}{9\sqrt3}=0.0260580$,
while with the uniform start no proper face wins
(Corollary~\ref{cor:simplex-modes}). Paper~C proves this window from its
general tie-window theorem. Dynamic programs over all bins at $N=90$, $300$ and
$1002$ confirm these modes at finite $N$ (numerically verified, see
Section~\ref{sec:verify}).
\end{remark}

\begin{remark}[the window at finite $N$]\label{rem:K3-finite}
The balanced bins of the edge cross the vertices exactly at Paper~A's root
$u_N$, for every $N\ge2$, since $\mu$ is constant on $\{1,2\}$. Indeed
$P^\mu_N(n_1,n_2,0)/P^\mu_N(N,0,0)=S_{(n_1,n_2)}(u)$, which is the polynomial of
Paper~A~\cite{paperA} at $u=r/p$. This crossing is at $T=Nu_N/(1+2u_N)$, which
differs from Paper~A's $z_N=Nu_N/(1+u_N)$ by $O(L^2/N)$. Near $t=a_2$ the lines
of the full face and of the edges through $3$ stay below those of the vertices,
so for large $N$ this crossing is the lower end of the window. For the 2 ends at
finite $N$ we use $z=Nu=T/p$. A heuristic calculation, which we do not prove,
gives
\[
 z-L+\tfrac12\log z=a_2+\frac1{8z}+\frac{z^2}{2N}+\cdots,\qquad
 z-L+\tfrac12\log z=t_{\rm ef}+\frac1{12z}+\frac{3z^2}{2N}+\cdots,
\]
where only $a_2+\frac1{8z}$ is proved, by Theorem~\ref{thm:simplex-crossing}.
The terms $z^2/N$ come from a discrete term $-\frac{(k-1)^2}2\frac{z^2}N$ in the
logarithm of~\eqref{eq:simplex-uniform-ratio}, and the upper end also uses a gain
$\frac1{24z}$ of the best full-face bin over the balanced one. Without the terms
$z^2/N$ the width is $0.0260580-0.0546957/L+\cdots$. Exact computations give
the widths $0.0843$, $0.0324$, $0.0223$ and $0.0214$ at $N=300$, $3000$,
$30000$ and $300000$, and the prediction gives $0.0802$, $0.0321$, $0.0226$ and
$0.0217$ (numerically verified, see Section~\ref{sec:verify}). Below $N$ of
about $10^4$ the terms $z^2/N$ are larger than the terms $1/L$.
\end{remark}
